%=====================================================================
\section{Introduction}
\label{sec:intro}

Network intrusion detection requires the assignment of individual traffic
records to attack categories. Manual inspection does not scale to the volume
of traffic carried by a modern network, so an \gls{ids} relies on automated
classification of traffic records. Supervised learning provides the
mechanism, in that a classifier learns the association between measured
flow characteristics and known attack categories.

A preceding study compared two single classifiers on an edited version of the
UNSW-NB15 dataset, described by Moustafa and Slay~\cite{moustafa2015}. A
classification tree outperformed \gls{knn} classification on every
performance measure that accounts for class imbalance. The tree nevertheless
detected the backdoor and analysis categories in fewer than half of the
cases, and the variation of a single tree under perturbation of the training
data remained unexamined. Single classification trees are known to be unstable,
in the sense that small changes to the training set produce large changes to
the induced tree~\cite{breiman1996}.

The present study investigates ensemble learning as a remedy. A \gls{rf},
introduced by Breiman~\cite{breiman2001}, was selected as the ensemble
method, and was compared against revised implementations of the
classification tree and the \gls{knn} model. The \gls{rf} shares the base
learner of the classification tree. Any difference in performance between
the two models is therefore attributable to the ensemble mechanism, rather
than to a change of hypothesis space. A bagged ensemble of trees was
evaluated as an intermediate model, so that the contribution of bootstrap
aggregation was separated from the contribution of random feature selection.

The tuning of the \gls{rf} was designed to answer five research questions,
fixed before any forest was trained:
\begin{enumerate}
\item Does the \gls{rf} improve on the tuned classification tree, and on
which attack categories does the improvement occur?
\item How does performance change with the number of trees, and does the
ensemble overfit as trees are added?
\item How does the size of the random feature subset trade the strength of
individual trees against the correlation between trees?
\item Do the depth and leaf size limits required by the single tree remain
beneficial inside an ensemble?
\item Does instance weighting in inverse proportion to class frequency remain
beneficial once predictions are averaged over many trees?
\end{enumerate}

The preceding study fitted selected pre-processing steps on the full
dataset and retained no held-out data. The present study removes both
limitations: every fitted transformation was learned inside the
cross-validation folds, and every model was evaluated once on a
generalisation set that played no part in any decision.

The remainder of the report is organised as follows.
Section~\ref{sec:background} presents the background on the three
classification approaches, on ensemble learning, and on the evaluation of
classifiers. Section~\ref{sec:data} describes the dataset and the
pre-processing. Section~\ref{sec:impl} describes the implementation.
Section~\ref{sec:empirical} describes the empirical process.
Section~\ref{sec:results} presents and discusses the results, and
Section~\ref{sec:conclusion} concludes the report.

%=====================================================================
\section{Background}
\label{sec:background}

The section provides the conceptual background required for the remainder of
the report. The section first describes the two single classifiers of the
preceding study, then presents ensemble learning in general and the
\gls{rf} in particular, and finally describes the measures and statistical
tests used to compare classifiers.

\subsection{Nearest neighbour classification}
\label{sec:bg-knn}

The \gls{knn} algorithm, analysed by Cover and Hart~\cite{cover1967}, is an
instance-based classifier. Training consists of the storage of the training
set. Classification of an unseen instance proceeds by the retrieval of the
$k$ stored instances closest to the unseen instance, followed by a vote among
the retrieved instances. For two instances represented by the feature vectors
$\mathbf{x}$ and $\mathbf{z}$ over $p$ features, the Euclidean distance is
\begin{equation}
\delta(\mathbf{x}, \mathbf{z}) = \sqrt{\sum_{j=1}^{p} (x_j - z_j)^2},
\label{eq:euclid}
\end{equation}
where $x_j$ denotes the value of feature $j$ of instance $\mathbf{x}$.
Equation~(\ref{eq:euclid}) sums squared differences of raw feature values,
so a feature of large numeric range dominates the distance irrespective of
the usefulness of the feature for discrimination. Feature scaling is
therefore a prerequisite of the algorithm.

Two control parameters govern the behaviour of the algorithm. A small value
of $k$ produces a flexible decision boundary that is sensitive to noise,
whereas a large value of $k$ smooths the boundary and biases prediction
toward the majority class. Under uniform weighting each retrieved neighbour
casts one vote. Under distance weighting each vote is scaled by the
reciprocal of $\delta$, so that closer neighbours exert greater influence.
Aggarwal \textit{et al.}~\cite{aggarwal2001} showed that the contrast
between the nearest and the furthest neighbour diminishes as the number of
dimensions grows, which penalises representations with many sparse
features.

\subsection{Classification trees}
\label{sec:bg-tree}

A classification tree, described by Breiman
\textit{et al.}~\cite{breiman1984}, recursively partitions the feature space.
Each internal node tests one feature against a threshold, and each leaf
assigns a class label. Construction is greedy: at each node the algorithm
selects the feature and threshold that maximise the reduction of an impurity
measure. For a node at which $\pi_j$ denotes the proportion of instances of
class $j$, and $c$ denotes the number of classes, the Gini index and the
entropy are
\begin{equation}
G = 1 - \sum_{j=1}^{c} \pi_j^2
\quad\text{and}\quad
H = -\sum_{j=1}^{c} \pi_j \log_2 \pi_j.
\label{eq:impurity}
\end{equation}
Entropy was employed by Quinlan~\cite{quinlan1986} as the split criterion of an early tree induction algorithm.

A tree grown until every leaf is pure fits the training set exactly and
generalises poorly. Pre-pruning limits growth through a maximum depth and a
minimum number of instances per leaf. Post-pruning grows a full tree and then
removes subtrees. Cost-complexity pruning~\cite{breiman1984} selects the
subtree of the full tree $h$ that minimises
\begin{equation}
R_\alpha(h) = R(h) + \alpha |h|,
\label{eq:ccp}
\end{equation}
where $R(h)$ denotes the weighted training misclassification cost of $h$,
$|h|$ denotes the number of leaves of $h$, and $\alpha \geq 0$ denotes the
complexity penalty. Every removed subtree is replaced by a single leaf that
predicts the majority class of the training instances that reach the
removed subtree. Increasing values of $\alpha$ therefore produce
progressively smaller trees.

Class imbalance is addressed through instance weights. Under weighting in
inverse proportion to class frequency, an instance of class $j$ receives the
weight $w_j = n / (c\,n_j)$, where $n$ denotes the number of training
instances and $n_j$ denotes the number of training instances of class $j$.
The weights enter the class proportions of Equation~(\ref{eq:impurity}), so
that a split which isolates a rare class achieves a larger impurity
reduction~\cite{he2009}.

\subsection{Ensemble learning}
\label{sec:bg-ensemble}

An ensemble combines the predictions of several base classifiers into a
single prediction. Dietterich~\cite{dietterich2000} identified three reasons
for the success of ensembles. The statistical reason is that a finite
training set supports many hypotheses of equal training accuracy, and an
average over the hypotheses reduces the risk of the selection of a poor
hypothesis. The computational reason is that greedy induction reaches a
local optimum, and different starting conditions reach different optima. The
representational reason is that a combination of hypotheses represents
decision boundaries that no single member represents.

The error of a learner decomposes into a bias term and a variance
term~\cite{geman1992}. Fully grown trees have low bias and high variance,
which makes the trees ideal members of an averaging ensemble: the average
retains the low bias and reduces the variance. An average reduces variance
only to the extent that the members disagree. Kuncheva and
Whitaker~\cite{kuncheva2003} reviewed measures of disagreement between
members, termed diversity, and examined the relationship between diversity
and ensemble accuracy.

Three families of ensemble methods dominate the literature. Bagging and
random forests train members independently on perturbed versions of the
training set and average the predictions. Boosting, of which AdaBoost by
Freund and Schapire~\cite{freund1997} and gradient boosting by
Friedman~\cite{friedman2001} are the best known instances, trains members
sequentially, so that each member concentrates on the errors of the
preceding members. Stacking, introduced by Wolpert~\cite{wolpert1992}, trains
a second-level learner on the outputs of heterogeneous members.

\subsection{Bagging}
\label{sec:bg-bagging}

Bagging, introduced by Breiman~\cite{breiman1996}, trains $\tau$ members on
$\tau$ bootstrap samples of the training set. A bootstrap sample is drawn
with replacement and has the size of the training set, so each sample omits
approximately $36.8\,\%$ of the distinct training instances. The members
vote, or average the class probabilities, to produce the prediction of the
ensemble. Breiman~\cite{breiman1996} showed that bagging improves unstable
learners, such as classification trees, and leaves stable learners, such as
the \gls{knn} algorithm, largely unchanged.

\subsection{Random forests}
\label{sec:bg-rf}

The \gls{rf}~\cite{breiman2001} extends bagging with a second source of
randomisation. At every node of every tree, the best split is sought among a
random subset of $q$ of the $p$ features, rather than among all features.
The idea follows the random subspace method of Ho~\cite{ho1998}. A value of
$q$ equal to $p$ reduces the \gls{rf} to bagging.

Breiman~\cite{breiman2001} bounded the generalisation error $\epsilon$ of a
forest by
\begin{equation}
\epsilon \leq \frac{\bar{\rho}\,(1 - s^2)}{s^2},
\label{eq:bound}
\end{equation}
where $s$ denotes the strength of the individual trees, defined as the
expected margin of a tree, and $\bar{\rho}$ denotes the mean correlation
between trees. Equation~(\ref{eq:bound}) identifies $q$ as the central
control parameter of the method. A small value of $q$ decorrelates the
trees, which reduces $\bar{\rho}$, but also weakens each tree, which reduces
$s$. The best value of $q$ balances the two effects and depends on the
dataset~\cite{probst2019}.

Breiman~\cite{breiman2001} additionally established, through the strong law
of large numbers, that the generalisation error converges as $\tau$ grows. An
increase of $\tau$ therefore cannot cause overfitting, and the choice of
$\tau$ trades computational cost against the residual variance of the
ensemble. Probst and Boulesteix~\cite{probst2018} showed that the performance
curve over $\tau$ is not always monotone for measures other than the error
rate, which motivates an empirical examination of the curve rather than an
assumption.

The trees of an \gls{rf} are conventionally grown to purity, because the
average removes the variance that pruning would otherwise control. Probst
\textit{et al.}~\cite{probst2019} reported that a limit on the leaf size
occasionally improves performance, particularly on noisy data. Class
weighting applies to the forest as to a single tree, either with weights
computed once from the full training set or with weights recomputed for
every bootstrap sample.

The \gls{rf} also offers a measure of feature importance, computed as the
total weighted impurity reduction attributable to each feature. Strobl
\textit{et al.}~\cite{strobl2007} showed that the impurity-based measure
favours features with many candidate split points. Permutation importance,
defined as the reduction of a performance measure after the random
permutation of the values of one feature in a held-out set, avoids the bias.

\subsection{Performance measures}
\label{sec:bg-measures}

Let $\mathbf{C}$ denote the confusion matrix of a classifier on a test set,
where the element $C_{jl}$ counts the instances of class $j$ predicted as
class $l$. The precision $a_j$, the recall $r_j$ and the $F_1$ measure
$F_j$ of class $j$ are
\begin{equation}
a_j = \frac{C_{jj}}{\sum_{l=1}^{c} C_{lj}}, \quad
r_j = \frac{C_{jj}}{\sum_{l=1}^{c} C_{jl}}, \quad
F_j = \frac{2\,a_j r_j}{a_j + r_j}.
\label{eq:prf}
\end{equation}
The macro-averaged $F_1$ measure $\bar{F}$, the balanced accuracy $\bar{r}$
and the accuracy $\eta$ are
\begin{equation}
\bar{F} = \frac{1}{c}\sum_{j=1}^{c} F_j, \quad
\bar{r} = \frac{1}{c}\sum_{j=1}^{c} r_j, \quad
\eta = \frac{\sum_{j=1}^{c} C_{jj}}{\sum_{j=1}^{c}\sum_{l=1}^{c} C_{jl}}.
\label{eq:macro}
\end{equation}
The accuracy $\eta$ weights every instance equally, and is therefore
dominated by the largest classes. The measures $\bar{F}$ and $\bar{r}$
weight every class equally, which suits problems in which the detection of
rare classes is the objective~\cite{sokolova2009}.

\subsection{Statistical comparison of classifiers}
\label{sec:bg-stats}

Dietterich~\cite{dietterich1998} recommended the McNemar
test~\cite{mcnemar1947} for the comparison of two classifiers evaluated on a
single test set. Let $u$ denote the number of test instances classified
correctly by the first classifier only, and $v$ the number classified
correctly by the second classifier only. Under the null hypothesis of equal
error rates, the statistic
\begin{equation}
\chi^2 = \frac{\left(|u - v| - 1\right)^2}{u + v}
\label{eq:mcnemar}
\end{equation}
follows a chi-squared distribution with one degree of freedom. The McNemar
test addresses accuracy only. A difference in $\bar{F}$ is instead assessed
with the paired bootstrap~\cite{efron1993}: test instances are resampled
with replacement, the difference in $\bar{F}$ between the two classifiers is
recomputed on every resample, and the percentiles of the resulting
distribution yield a confidence interval. When several pairwise tests are
performed, the procedure of Holm~\cite{holm1979} adjusts the $p$-values so
that the probability of one or more false rejections across the family
remains below the significance level.

%=====================================================================
\section{Data and pre-processing}
\label{sec:data}

The section describes the dataset, the properties of the dataset that bear on
the choice and configuration of the models, and the pre-processing applied
to the input of each model. The section first reports the data quality
analysis, then describes the corrections shared by all models, and finally
describes the model-specific pre-processing candidates that were evaluated.

\subsection{Dataset}
\label{sec:dataset}

The supplied file contains \num{257673} network flow records, each described
by a record identifier, 42 descriptive features and a target feature. The
target feature assigns every record to one of ten categories, namely normal
traffic and nine attack categories. The descriptive features comprise three
nominal features, namely the transaction protocol, the connection state and
the application service, and 39 numeric features that describe packet counts,
byte counts, timing, and connection statistics~\cite{moustafa2016}.

The data quality analysis of the preceding study established eight defects.
The record identifier is unique and ordered by category. The service feature
contains \num{141321} missing values. The binary feature that indicates a
\gls{ftp} login contains 26 values greater than one. Numeric ranges span nine
orders of magnitude, 24 feature pairs correlate above 0.85, and the protocol
feature has 133 levels. Finally, \num{94928} records, or $36.8\,\%$ of the
file, are exact duplicates once the identifier is disregarded. After the
removal of the duplicates, \num{162745} records remain, of which normal
traffic accounts for $52.7\,\%$ and the worms category for $0.1\,\%$.

The present study established a ninth property of the data, which the
preceding study did not examine. After deduplication, \num{11506} records,
or $7.1\,\%$ of the dataset, share identical values on all 42 descriptive
features with at least one record of a different category. No classifier
that receives only the 42 features distinguishes such records. The conflicts
are concentrated in specific categories: $81.4\,\%$ of the backdoor records,
$77.9\,\%$ of the analysis records and $34.1\,\%$ of the denial of service
records belong to a group of identical feature vectors that spans several
categories, against $0.5\,\%$ of the normal records and $0.7\,\%$ of the
shellcode records. The largest such groups mix the reconnaissance,
backdoor, denial of service, exploits, analysis and fuzzers categories.
The property bounds the performance attainable on the affected categories,
because a group of identical feature vectors receives a single prediction.

\subsection{Corrections shared by all models}
\label{sec:shared}

Four corrections were applied identically to the input of every model. The
record identifier was removed, because the identifier carries no
generalisable information and the ordering of the file associates identifier
ranges with categories. Missing service values were encoded as an additional
level, because the absence of a notable application service is itself an
observation~\cite{moustafa2015}. Values of the \gls{ftp} login indicator greater
than one were clipped to one, because any such value records the occurrence
of a login. Duplicate records were removed, because a duplicate placed in
both a training partition and a test partition inflates the estimated
performance.

None of the four corrections learns from the data, so the corrections were
applied once, before the partitioning of the data. Every remaining
transformation learns a quantity from the data, such as a median, a level
vocabulary or a correlation. Every such transformation was fitted on the
training partition of each fold and then applied, unchanged, to the
corresponding test partition.

\subsection{Pre-processing candidates for nearest neighbour classification}
\label{sec:prep-knn}

The \gls{knn} input requires three decisions, and the preceding study left
two of the decisions untested. The first decision is the scaling method.
Four candidates were evaluated: standardisation to zero mean and unit
variance, scaling by the median and the interquartile range, scaling to the
range $[0, 1]$, and the absence of scaling. Interquartile scaling was
expected to perform best, because the median and the interquartile range are
insensitive to the extreme values present in the data.

The second decision is the encoding of nominal features. Every level was
encoded as a binary indicator feature, because an integer code imposes an
artificial order that Equation~(\ref{eq:euclid}) would treat as a meaningful
distance. The 133 protocol levels then produce 133 sparse indicators. Two
candidates were evaluated: retention of every level, and the retention of
the ten most frequent protocol levels and the five most frequent state levels
with the remaining levels merged into a single level. The merge was expected
to help modestly, through the mechanism described by Aggarwal
\textit{et al.}~\cite{aggarwal2001}.

The third decision concerns two optional transformations, namely a
logarithmic transformation of non-negative features with a skewness above
one, and the removal of one feature from every pair with an absolute
correlation above 0.95. Each transformation was evaluated as an addition to
the best configuration of the first two decisions. The preceding study
found that both transformations degraded performance, but the finding rested
on transformations fitted to the full dataset.

\subsection{Pre-processing candidates for tree-based models}
\label{sec:prep-tree}

Scaling, logarithmic transformation and the removal of correlated features
were not applied to the input of the classification tree or the tree
ensembles. A threshold test depends only on the order of values within a
feature, so a monotone transformation leaves every candidate split
unchanged. A redundant feature yields no additional impurity reduction, so
the tree removes redundancy implicitly.

The encoding of nominal features does affect tree-based models, and two
encodings were evaluated for each tree-based model. Integer encoding assigns
one code per level, which keeps the number of features at 42 but forces the
tree to isolate a level through several threshold tests. Indicator encoding
allows a single test to isolate one level, at the cost of an increase to 193
features. For the classification tree, both encodings were
additionally evaluated with and without the merge of rare levels.

%=====================================================================
\section{Implementation}
\label{sec:impl}

The section describes the software, the implementation of each model, and the
measures taken for reproducibility and computational tractability.

\subsection{Software}

The study was implemented in Python, version 3.13. Models, resampling,
scaling and performance measures were provided by scikit-learn, version
1.9~\cite{pedregosa2011}, through the classes
\texttt{DecisionTreeClassifier}, \texttt{RandomForestClassifier},
\texttt{NearestNeighbors}, \texttt{StratifiedKFold}, \texttt{RobustScaler},
\texttt{StandardScaler} and \texttt{MinMaxScaler}. Data manipulation was
performed with pandas, version 3.0, and numeric computation with NumPy,
version 2.5. The source code is publicly
available.\footnote{\url{https://github.com/J-Burden/26905167-rw741-assignment3}}

The level merge and both encodings were implemented as transformer classes
with separate fitting and application methods, so that the level vocabulary
of each fold was learned from the training partition only.

\subsection{Nearest neighbour search}

The \gls{knn} model was implemented through exhaustive search rather than
through a spatial index. A spatial index, such as a $k$-d tree, partitions the
feature space to avoid the computation of every distance, but degenerates to
exhaustive search in high-dimensional spaces, and the indicator encoding
produces between 62 and 193 features. A tuning sweep over several values of
$k$ and both weighting schemes requires only one search per fold, because
the neighbours returned by an exhaustive search are ordered by distance. The
list of the $k$ nearest neighbours is therefore a prefix of the list for
any larger value of $k$. Algorithm~\ref{alg:reuse} states the procedure.

\begin{algorithm}[!t]
\caption{Neighbour reuse for a control parameter sweep}
\label{alg:reuse}
\begin{algorithmic}[1]
\For{each fold}
  \State Fit the pre-processing on the training partition
  \State Apply the pre-processing to both partitions
  \State Retrieve the largest candidate number of nearest training instances
         of every test instance, in order of increasing distance
  \For{each candidate value of $k$}
    \State Restrict every retrieved list to the first $k$ entries
    \State Predict the class of greatest vote count under uniform weighting
    \State Predict the class of greatest total reciprocal distance under
           distance weighting
  \EndFor
\EndFor
\end{algorithmic}
\end{algorithm}

Where a neighbour lay at distance zero, every neighbour at distance zero
received a weight of one and every other neighbour a weight of zero, which
matches the convention of scikit-learn.

\subsection{Tree-based models}

The classification tree was implemented through
\texttt{DecisionTreeClassifier}, which grows binary trees and supports both
impurity measures of Equation~(\ref{eq:impurity}), class weighting, and the
cost-complexity pruning of Equation~(\ref{eq:ccp}). The implementation
requires numeric input without missing values, which the encodings of
Section~\ref{sec:prep-tree} provide. The bagged ensemble and the \gls{rf}
were implemented through \texttt{RandomForestClassifier}, with $q$ set to
$p$ for bagging. The ensemble predicts the class of greatest mean class
probability over the trees, rather than the class of greatest vote count.

The examination of performance as a function of $\tau$ did not require a
separate forest for every value of $\tau$. Each forest member is trained
independently of the other members, so the first $t$ members of a forest of
$\tau$ trees form a valid forest of $t$ trees. A single forest of 300 trees
was therefore trained per fold. The class probabilities of the members were
accumulated in order, and the prediction of the forest was recorded after
the addition of every member.

The diversity of every forest was measured on the test partition of each
fold. The strength of the members was measured as the mean $\bar{F}$ of the
individual trees. The correlation between members was measured as the mean
pairwise agreement, defined as the proportion of instances for which two
trees predict the same class~\cite{kuncheva2003}, over every pair among 25
randomly selected trees.

\subsection{Reproducibility and computational budget}

Fixed seeds were supplied to the partitioning, to the folds, and to every
tree and forest, so that every model was evaluated on identical partitions.
The experiments were constrained to two processor cores, which limited the
size of the tuning grids and motivated the use of 100 trees during the
tuning of the \gls{rf}.
